%% kapa: exam=1 round=36 year=2025 period=1 form=A subject=civil_law subject_ko=민법 seq=18 number=18
%% source: Q-Net 감정평가사 기출문제 (cst00309) artlSeq=5231525 file=raw/1008/5231525/2025년 제36회 감정평가사 1차시험 1교시 문제지 원본.pdf
%% conversion: kordoc 4.24.51 -> tools/gen1.py (auto); review=proofread-sonnet; flags=none
\begin{question}{18}
소멸시효에 관한 설명으로 옳지 않은 것은? (다툼이 있으면 판례에 따름)
\begin{choices}
  \item 소유권에 기한 물권적 청구권은 소멸시효의 대상이 되지 않는다.
  \item 정지조건부 권리는 조건 미성취의 동안에도 소멸시효가 진행된다.
  \item 3년의 단기소멸시효가 적용되는 ‘1년 이내의 기간으로 정한 채권’이란 1년 이내의 정기로 지급되는 채권을 말한다.
  \item 채권의 이행기가 도래한 후 채권자와 채무자가 이행기를 유예하기로 합의한 경우, 소멸시효는 변경된 이행기가 도래한 때부터 다시 진행한다.
  \item 변제기가 도래하여 지급명령에서 확정된 채권은 단기의 소멸시효에 해당하는 것이라도 그 소멸시효기간이 10년으로 연장된다.
\end{choices}
\end{question}
\answer{2}
